\documentclass[11pt,a4paper]{article}

\usepackage[margin=2.5cm]{geometry}
\usepackage{amsmath,amssymb,amsthm}
\usepackage{booktabs}
\usepackage{enumitem}
\usepackage{algorithm}
\usepackage[noend]{algpseudocode}
\usepackage[hidelinks]{hyperref}

\newtheorem{theorem}{Theorem}[section]
\newtheorem{proposition}[theorem]{Proposition}
\newtheorem{lemma}[theorem]{Lemma}
\newtheorem{corollary}[theorem]{Corollary}
\theoremstyle{definition}
\newtheorem{definition}[theorem]{Definition}
\newtheorem{example}[theorem]{Example}
\theoremstyle{remark}
\newtheorem{remark}[theorem]{Remark}

\newcommand{\proc}[1]{\textnormal{\textsc{#1}}}
\newcommand{\N}{\mathbb{N}}
\newcommand{\Ops}{\operatorname{Ops}}
\newcommand{\val}{\operatorname{val}}
\newcommand{\leaves}{\operatorname{lv}}
\newcommand{\mask}{\mu}
\newcommand{\str}[1]{\texttt{#1}}

\title{The Countdown Numbers Problem:\\ a Canonical Formulation}
\author{}
\date{October 2026}

\begin{document}
\maketitle

\begin{abstract}
In the numbers round of the television game \emph{Countdown}, contestants combine six dealt numbers with
the four arithmetic operations to reach a three-digit target. We give a precise mathematical definition
of the problem solved by the implementations in the accompanying repository~\cite{repo}. The definition goes beyond ``reach the
target'': it fixes the rules for intermediate values, an objective that prefers exact and then
simpler answers, and a total tie-breaking order that makes the answer, including its printed text,
a function of the input. We prove that the definition is well posed, show that it agrees with the
staged ``exact first, otherwise closest'' procedure the solvers carry out, and record the structural facts
the fast solvers rely on. We then describe the algorithms the solvers use and prove them correct, from the
sequential reference to massively parallel GPU versions, and close with statistics and timings for every implementation on a 10{,}000-round benchmark set.
\end{abstract}

\section{Introduction}

A Countdown numbers round deals a \emph{pool} of six positive integers and a \emph{target} between
100 and 999. Using each number at most once, together with addition, subtraction, multiplication and
division, players try to make the target. Every intermediate result must be a positive integer, and not
every number has to be used. If the target cannot be reached, the closest attainable value wins.

The repository~\cite{repo}, at \url{https://github.com/ZippyManaic/countdown}, contains solvers for this game in more than thirty languages, from x86-64 assembly and line-numbered
BASIC to CUDA, plus seven versions that share each round across several threads. They are required to produce \emph{byte-identical} output, so the problem they
solve has to be defined down to which of several equally good expressions is printed and how it is
bracketed. This paper gives that definition. Section~\ref{sec:instances} describes the inputs.
Section~\ref{sec:expressions} defines legal expressions and the values they reach.
Section~\ref{sec:objective} defines the objective and Section~\ref{sec:canonical} the canonical answer.
Section~\ref{sec:structure} collects the properties that efficient solvers exploit,
Section~\ref{sec:algorithms} gives the algorithms with their correctness proofs and
Section~\ref{sec:data} reports measurements of the problem and Section~\ref{sec:performance} the
timings of every implementation. Section~\ref{sec:optimisation} compares each language's straight port with
its optimised version, and Section~\ref{sec:parallel} shares a single round across CPU threads. Hutton~\cite{hutton} gave an earlier formalisation that
lists every solution. Here the task is to choose exactly one.

\section{Instances}\label{sec:instances}

\begin{definition}[Instance]
An \emph{instance} is a pair $(t, p)$ in which the \emph{target} $t \in \N_{\ge 1}$ is a positive integer
and the \emph{pool} $p = (p_0, \dots, p_{n-1})$ is a finite sequence of $n \ge 1$ positive integers.
Write $[n] = \{0, \dots, n-1\}$ for the set of \emph{positions}.
\end{definition}

The pool is a \emph{sequence} rather than a multiset. Positions distinguish equal numbers, and
Section~\ref{sec:canonical} shows that the order of the pool can change the printed answer. It cannot change
how good that answer is.

\begin{definition}[Subsets and masks]
For $S \subseteq [n]$ let $\mask(S) = \sum_{i \in S} 2^i$ be its \emph{mask}. The map $\mask$ is a bijection
from subsets of $[n]$ to $\{0, \dots, 2^n - 1\}$, and we freely identify a subset with its mask. Write
$|S|$ for the number of positions in $S$.
\end{definition}

\paragraph{The game's distribution.}
A round of the game is dealt as follows. The target is uniform on $\{100, \dots, 999\}$. The number of
\emph{large} numbers $\ell$ is uniform on $\{1, 2\}$. The $6 - \ell$ \emph{small} numbers are drawn without
replacement from a deck holding two cards of each of $1, \dots, 10$, and the $\ell$ large numbers are drawn
without replacement from $\{25, 50, 75, 100\}$. The pool lists the small numbers followed by the large ones.
The solvers accept any pool with $1 \le n \le 6$, and the definitions below hold for every $n$.

\begin{proposition}\label{prop:pools}
Up to order, the game can deal exactly $9{,}498$ distinct pools, so it has $9{,}498 \times 900 = 8{,}548{,}200$
distinct instances.
\end{proposition}

\begin{proof}
A multiset of $k$ small numbers using each of $1, \dots, 10$ at most twice corresponds to the coefficient of
$x^k$ in $(1 + x + x^2)^{10}$. That coefficient is $615$ for $k = 4$ and $1{,}452$ for $k = 5$. Two large numbers
can be chosen in $\binom{4}{2} = 6$ ways and one in $4$ ways. The large and small numbers never coincide,
so the count is $6 \cdot 615 + 4 \cdot 1{,}452 = 9{,}498$.
\end{proof}

\section{Expressions and reachable values}\label{sec:expressions}

\begin{definition}[Legal operations]\label{def:ops}
For integers $x \ge y \ge 1$, the \emph{legal operations} on the ordered pair $(x, y)$ are listed below in
their \emph{operator order}. Each one is legal only when its condition holds.
\[
\begin{array}{llll}
\text{rank} & \text{operator} & \text{result} & \text{condition} \\ \midrule
0 & + & x + y & \text{always} \\
1 & - & x - y & x > y \\
2 & \times & x \cdot y & y \ne 1 \\
3 & \div & x / y & y \ne 1 \text{ and } y \mid x
\end{array}
\]
Write $\Ops(x, y)$ for the list of legal pairs $(r, \circ)$ of result and operator. Every result is a positive
integer.
\end{definition}

Operations always take the larger operand first, so subtraction never goes negative and division never
produces a fraction below one. The game's rules additionally allow $x - x = 0$, $x \times 1$ and $x \div 1$.
Proposition~\ref{prop:restriction} shows that excluding them never changes how good the best answer is.

\begin{definition}[Expression]\label{def:expr}
An \emph{expression} over $p$ is a finite binary tree $T$ with a value $\val(T)$ and a leaf set
$\leaves(T) \subseteq [n]$, built by two rules:
\begin{enumerate}[label=(\roman*)]
\item for each $i \in [n]$ the leaf $\langle i \rangle$ is an expression with $\val = p_i$ and $\leaves = \{i\}$;
\item if $T_1$ and $T_2$ are expressions with $\leaves(T_1) \cap \leaves(T_2) = \emptyset$ and
      $\val(T_1) \ge \val(T_2)$, and $(r, \circ) \in \Ops(\val(T_1), \val(T_2))$, then the node
      $(T_1 \circ T_2)$ is an expression with $\val = r$ and $\leaves = \leaves(T_1) \cup \leaves(T_2)$.
\end{enumerate}
\end{definition}

Disjointness of leaf sets says that each dealt number is used at most once. The number of numbers an
expression \emph{uses} is $|\leaves(T)|$.

\begin{definition}[Reachable values]
For non-empty $S \subseteq [n]$, let
\[
V(S) = \{\val(T) : T \text{ an expression with } \leaves(T) = S\}.
\]
\end{definition}

A \emph{split} of $S$ is an unordered pair $\{L, R\}$ of non-empty sets with $L \cup R = S$ and
$L \cap R = \emptyset$. Because $\mask$ is injective, each split has a unique representative with
$\mask(L) < \mask(R)$.

\begin{theorem}[Subset recursion]\label{thm:recursion}
$V(\{i\}) = \{p_i\}$, and for $|S| \ge 2$
\[
V(S) = \bigcup_{\{L, R\} \text{ split of } S} \;
       \bigcup_{a \in V(L),\, b \in V(R)}
       \{\, r : (r, \circ) \in \Ops(\max(a, b), \min(a, b)) \,\}.
\]
\end{theorem}

\begin{proof}
A tree whose leaf set is a singleton must be that leaf, which gives the base case. For $|S| \ge 2$ a tree $T$
with $\leaves(T) = S$ is a node $(T_1 \circ T_2)$. Its children's leaf sets form a split of $S$, and their
values lie in $V(L)$ and $V(R)$, so $\val(T)$ belongs to the right-hand side. Conversely, take $a \in V(L)$
and $b \in V(R)$ for a split, witnessed by trees $A$ and $B$. Put the one with the larger value first, taking
$A$ when $a = b$. Then each legal operation builds a node whose value is the corresponding result, and its
leaf set is $S$.
\end{proof}

\begin{proposition}[Value bound]\label{prop:bound}
Every $v \in V(S)$ satisfies $1 \le v \le \prod_{i \in S} (p_i + 1) - 1$.
\end{proposition}

\begin{proof}
Use induction on $|S|$, where the base case is $p_i \le (p_i + 1) - 1$. For $x \ge y \ge 1$ each legal result is
at most $\max(x + y, x y) \le x y + x + y = (x + 1)(y + 1) - 1$, and the bounds multiply over a split.
\end{proof}

For a game pool this is below $101 \cdot 76 \cdot 11^2 \cdot 10^2 < 2^{27}$, so every value fits in 32 bits.
Several solvers size a bitset of seen values by this bound.

\begin{proposition}[Restricted rules lose nothing]\label{prop:restriction}
Allow also the results $x - x = 0$, $x \times 1$ and $x \div 1$, and let $d^\dagger$ and $k^\dagger$ be the
optimal distance and number count of Definition~\ref{def:objective} under these looser rules. If $t \ge 1$
then the restricted rules attain the same $d^\dagger$ with the same $k^\dagger$.
\end{proposition}

\begin{proof}
It suffices to show that any loose expression using one of the extra moves can be replaced by one that is
no worse in distance and strictly better in distance or count. A node $T_1 \times 1$ or $T_1 \div 1$ can be
replaced by $T_1$, which has the same value and one leaf fewer, and the replacement propagates up
unchanged. Next consider a zero. Its parent combines it with some $w \ge 0$ through $w + 0$, $w - 0$ or
$w \times 0$, giving $w$ or $0$. In the first case the zero subtree can be deleted, which keeps the value and
uses fewer leaves. Otherwise zero propagates to the root, so the whole expression has distance $t$. Take a
lowest zero node $x - x$ with $x \ge 1$. If $x \ne 1$, replace it by $x \div x = 1$, which uses the same
leaves. If $x = 1$, take either operand subtree, which has value $1$ and fewer leaves. Either way a value of
$1$ is reached, at distance $t - 1 < t$.
\end{proof}

\section{The objective}\label{sec:objective}

\begin{definition}[Candidates and the objective]\label{def:objective}
The \emph{candidates} of an instance are the pairs $(S, v)$ with $\emptyset \ne S \subseteq [n]$ and
$v \in V(S)$. Each candidate has the key
\[
\kappa(S, v) = \bigl(\, |t - v|,\; |S|,\; \mask(S),\; v \,\bigr) \in \N^4,
\]
and the \emph{optimal candidate} $(S^*, v^*)$ is the candidate with the lexicographically least key.
Call $d^* = |t - v^*|$ the \emph{optimal distance} and $k^* = |S^*|$ the \emph{optimal count}.
\end{definition}

The four components, in priority order, say: get as close to the target as possible, then use as few
numbers as possible, then prefer the subset with the smallest mask, then the smaller value. The last
component only separates $t - d$ from $t + d$, and it picks the value below the target.

\begin{proposition}
The optimal candidate exists and is unique.
\end{proposition}

\begin{proof}
There is at least one candidate, since $(\{0\}, p_0)$ qualifies, and there are finitely many, since each $V(S)$
is finite by Proposition~\ref{prop:bound}. The key determines $S$ through $\mask$ and also $v$, so distinct
candidates have distinct keys.
\end{proof}

The solvers do not minimise $\kappa$ directly. They follow the staged procedure below, which mirrors how
the game is played.

\begin{proposition}[Staged form]\label{prop:staged}
Let $F = [n]$ be the whole pool. The optimal candidate is produced by the following procedure.
\begin{enumerate}[label=(\alph*)]
\item If $t \in V(S)$ for some $S$, let $k$ be the least such $|S|$. Return $(S, t)$ for the least mask $S$ of size
      $k$ with $t \in V(S)$.
\item Otherwise let $(S_1, v_1)$ minimise $\kappa$ over candidates with $S \ne F$, and let $v_2$ minimise
      $(|t - v|, v)$ over $v \in V(F)$. Return $(F, v_2)$ if $|t - v_2| < |t - v_1|$, and $(S_1, v_1)$ otherwise.
\end{enumerate}
\end{proposition}

\begin{proof}
Case (a) is the definition restricted to distance $0$. In case (b), $F$ is the only subset of size $n$, so a
candidate on $F$ has a count larger than every other candidate. A candidate on $F$ therefore beats
$(S_1, v_1)$ under $\kappa$ exactly when it is strictly closer. Among candidates on $F$, the mask and size
are fixed, so $\kappa$ reduces to $(|t - v|, v)$. When $n = 1$ the set of candidates with $S \neq F$ is empty
and the procedure returns the single number.
\end{proof}

Stage (b) only ever needs the closest values of $V(F)$, never the whole set. This is why solvers can stream
the full pool's results without storing them.

\section{The canonical answer}\label{sec:canonical}

Many expressions can share the optimal value and subset. For example $4 + 2$ and $2 \times 3$ might both
give $6$ from different subsets, and a single subset typically has many trees per value. To make the output
a function of the input we fix one tree, the first derivation in a prescribed enumeration order.

\begin{definition}[Split order]
For $|S| \ge 2$ the \emph{split order} $\Sigma(S)$ is the sequence of sets $L$ such that
$\emptyset \ne L \subsetneq S$ and $\mask(L) < \mask(S \setminus L)$, listed in \emph{decreasing} order of
$\mask(L)$. Each split $\{L, S \setminus L\}$ appears exactly once. Write $\sigma_S(L)$ for the position of $L$
in this sequence.
\end{definition}

This is the order produced by the standard sub-mask loop
$L \leftarrow (S - 1) \mathbin{\&} S$, then $L \leftarrow (L - 1) \mathbin{\&} S$ repeatedly, keeping only
$L < S \oplus L$.

\begin{definition}[Canonical derivation]\label{def:derivation}
Let $|S| \ge 2$ and $v \in V(S)$. A \emph{derivation} of $v$ from $S$ is a tuple $(L, a, b, \circ)$ with
$L \in \Sigma(S)$, $R = S \setminus L$, $a \in V(L)$, $b \in V(R)$ and $(v, \circ) \in \Ops(\max(a, b), \min(a, b))$.
The \emph{canonical derivation} $\omega(S, v)$ is the derivation that minimises
\[
\bigl(\, \sigma_S(L),\; a,\; b,\; \operatorname{rank}(\circ) \,\bigr)
\]
lexicographically, where the rank is the one in Definition~\ref{def:ops}.
\end{definition}

The canonical derivation exists because $v \in V(S)$ and Theorem~\ref{thm:recursion} guarantee at least one
derivation. It is unique because the key, together with $S$, determines the tuple. In implementation terms
it is the \emph{first} way of making $v$ found when the splits are visited in order, each half's values in
ascending order and the operators in operator order.

\begin{definition}[Canonical tree]
Define $\tau(S, v)$ for every candidate by recursion on $|S|$. If $S = \{i\}$, then $\tau(S, v) = \langle i \rangle$.
Otherwise let $\omega(S, v) = (L, a, b, \circ)$ and $R = S \setminus L$, and set
\[
\tau(S, v) =
\begin{cases}
\bigl(\tau(L, a) \circ \tau(R, b)\bigr) & \text{if } a \ge b,\\[2pt]
\bigl(\tau(R, b) \circ \tau(L, a)\bigr) & \text{if } a < b.
\end{cases}
\]
\end{definition}

The recursion terminates because $|L|, |R| < |S|$. When the operands are equal the half with the smaller mask
is written first. Every subtree of a canonical tree is itself canonical for its own subset and value. This
is what lets a solver store a single origin per reachable value.

\begin{definition}[Rendering]
The text $\rho(T)$ of an expression is defined by
$\rho(\langle i \rangle) = $ the decimal digits of $p_i$, and
\[
\rho\bigl((T_1 \circ T_2)\bigr) = \str{(}\;\rho(T_1)\;\str{\textvisiblespace}\;s(\circ)\;\str{\textvisiblespace}\;\rho(T_2)\;\str{)},
\]
where $s(+) = \str{+}$, $s(-) = \str{-}$, $s(\times) = \str{*}$ and $s(\div) = \str{/}$, and juxtaposition is
concatenation. The \emph{display text} $\hat\rho(T)$ removes the first and last characters of $\rho(T)$ when $T$
is a node and equals $\rho(T)$ when $T$ is a leaf.
\end{definition}

Every node is fully bracketed, with the larger operand on the left, so no precedence conventions are
needed to read the text back. Only the outermost brackets are dropped.

\begin{definition}[The Countdown problem]\label{def:problem}
Given an instance $(t, p)$, compute
\[
\operatorname{Solve}(t, p) = \bigl(\, v^*,\; \hat\rho(\tau(S^*, v^*)) \,\bigr),
\]
where $(S^*, v^*)$ is the optimal candidate of Definition~\ref{def:objective}.
\end{definition}

\begin{example}
For $t = 125$ and $p = (1, 5, 7, 3, 75, 50)$ no single number equals $125$, but $75 + 50$ does, so $k^* = 2$ and
$\operatorname{Solve} = (125, \str{75 + 50})$.
\end{example}

\begin{example}
For $t = 423$ and $p = (6, 2, 2, 1, 1, 25)$ the target needs all six numbers:
\[
\operatorname{Solve}(t, p) = \bigl(423,\; \str{(25 * ((6 * (2 + 1)) - 1)) - 2}\bigr).
\]
\end{example}

\begin{example}
For $t = 431$ and $p = (3, 2, 3, 4, 2, 25)$ the target is unreachable. Both $430$ and $432$ lie at distance $1$.
The key's value component would favour $430$, but $430$ needs all six numbers while $432$ needs five, so the
count decides:
$\operatorname{Solve}(t, p) = \bigl(432,\; \str{((25 * 3) - 3) * (4 + 2)}\bigr)$.
\end{example}

\begin{remark}[Invariance]\label{rem:invariance}
Permuting the pool permutes the positions and preserves every $V(S)$ up to relabelling, so $d^*$, $k^*$ and $v^*$
can change only through the mask and value components of the key. The two \emph{costs} $(d^*, k^*)$ are
invariant. The \emph{display text} is not, because masks and split orders depend on positions. This is why
the lookup-table solver, which answers for the sorted pool, matches the reference text on only 69\% of rounds
while never being worse.
\end{remark}

\begin{definition}[Equally good answers]
An answer $(v, e)$ to $(t, p)$ is \emph{equally good} as $\operatorname{Solve}(t, p)$ when $e$ is a legal
expression over $p$ with value $v$, $|t - v| = d^*$ and $e$ uses $k^*$ numbers.
\end{definition}

The benchmark checker counts an answer as \emph{identical} when it matches $\operatorname{Solve}$ byte for byte,
as a \emph{tie} when it is equally good, and as \emph{wrong} otherwise. In practice the checker approximates
legality by evaluating the text and comparing the value and the count of numbers.

\section{Structure used by the solvers}\label{sec:structure}

\begin{proposition}[Number of splits]
A set of size $k \ge 2$ has $2^{k-1} - 1$ splits, and the pool as a whole has
$\sum_{k} \binom{n}{k} (2^{k-1} - 1) = \tfrac12 (3^n - 2^{n+1} + 1)$ splits across all its subsets. For $n = 6$ this
is $301$.
\end{proposition}

\begin{proof}
Ordered pairs $(L, R)$ with both non-empty number $2^k - 2$, and each split counts twice. For the total, use
$\sum_k \binom{n}{k} 2^k = 3^n$.
\end{proof}

\begin{proposition}[Smallest subsets first]\label{prop:early}
If $t \in V(S)$ for some $S$ with $|S| = k$, then $k^* \le k$ and no subset of size greater than $k$ affects
the answer. In particular, a solver that builds $V$ in increasing order of size can stop at the first size
whose subsets contain the target. Among those subsets, the smallest mask is the answer.
\end{proposition}

\begin{proof}
Any candidate with $|S| > k$ has a key with first component at least $0$ and second component greater than
$k$, so it loses to $(S, t)$. The canonical tree of the answer only involves subsets of $S^*$.
\end{proof}

Testing whether $t \in V(S)$ does not require building $V(S)$. It is enough to look, for each split and each
value on one side, for the few values on the other side that combine to $t$.

\begin{lemma}[Partners]\label{lem:partners}
Let $a, b, t \ge 1$. Some legal operation on $(\max(a, b), \min(a, b))$ yields $t$ if and only if $b$ belongs to
the set of positive integers among
\[
t - a, \qquad a - t, \qquad a + t, \qquad \tfrac{t}{a}\ \ (a \mid t,\ \min(a, \tfrac{t}{a}) \ne 1), \qquad
\tfrac{a}{t}\ \ (t \mid a,\ \tfrac{a}{t} \ne 1), \qquad a t\ \ (a \ne 1).
\]
\end{lemma}

\begin{proof}
These are the solutions for $b$ of $a + b = t$, $|a - b| = t$ with $a \ne b$, $a b = t$ with neither factor $1$,
and the quotient of the larger by the smaller equalling $t$ with divisor not $1$. Positivity of $t$ gives
$a \ne b$ in the subtraction cases. For the quotient $a / b = t$ the divisor is $b = a / t \le a$. For
$b / a = t$ it is $a$, and $b = a t \ge a$.
\end{proof}

So $t \in V(S)$ can be decided with at most six membership queries per value on the smaller half of each split.
The canonical derivation is the least $(\sigma_S(L), a, b, \operatorname{rank})$ among the hits. The fast
solvers also use this lemma recursively, to decide $t \in V(S)$ for $|S| = n$ without first building
$V$ for subsets of size $n - 1$.

\begin{remark}[Hardness]
With only addition allowed, deciding whether $d^* = 0$ is the Subset Sum problem, which is NP-complete when
numbers are written in binary. The game avoids this by fixing $n = 6$ and keeping numbers small.
Every instance is then a constant-size computation, and the interest lies in its constant factors.
\end{remark}

\section{Algorithms}\label{sec:algorithms}

This section describes how the implementations compute $\operatorname{Solve}$ and proves that they compute
it exactly. Section~\ref{sec:reference} gives the reference algorithm that every byte-identical solver
follows. Sections~\ref{sec:cpu} and~\ref{sec:gpu} cover the refinements used on CPUs and GPUs.
Section~\ref{sec:approx} covers the two versions that return equally good rather than identical answers.

\subsection{The reference algorithm}\label{sec:reference}

The reference algorithm stores, for every subset $S$ it builds, the sorted list of $V(S)$ with one
\emph{origin} per value, recording the derivation that made it. Two procedures act on a subset.
$\proc{Build}$ enumerates every derivation in canonical order and keeps the first one for each value.
$\proc{Find}$ decides whether a value is in $V(S)$ without building $V(S)$, by scanning one half of
each split and looking up the partners of Lemma~\ref{lem:partners} in the other.

\begin{algorithm}[ht]
\caption{Reference solver for $\operatorname{Solve}(t, p)$}\label{alg:reference}
\begin{algorithmic}[1]
\Procedure{Build}{$S$}
  \State $W \gets$ empty map from values to derivations
  \For{$L \in \Sigma(S)$ in order, with $R = S \setminus L$}
    \For{$a \in V(L)$ ascending, $b \in V(R)$ ascending}
      \For{$(r, \circ) \in \Ops(\max(a, b), \min(a, b))$ in operator order}
        \If{$r \notin W$} $W[r] \gets (L, a, b, \circ)$ \EndIf
      \EndFor
    \EndFor
  \EndFor
  \State store $V(S) \gets$ the keys of $W$ sorted ascending, each with its derivation
\EndProcedure
\Statex
\Procedure{Find}{$S, v$}
  \For{$L \in \Sigma(S)$ in order, with $R = S \setminus L$}
    \State $H \gets \{(a, b, \circ) : a \in V(L),\ b \in V(R),\ (v, \circ) \in \Ops(\max(a, b), \min(a, b))\}$
    \State \Comment{enumerate $H$ by scanning the shorter of $V(L)$ and $V(R)$ and looking up partners}
    \If{$H \neq \emptyset$} \Return $(L, \min H)$ \EndIf
  \EndFor
  \State \Return $\bot$
\EndProcedure
\Statex
\Procedure{Solve}{$t, p$}
  \For{$i \in [n]$} store $V(\{i\}) \gets \{p_i\}$ \EndFor
  \If{$p_i = t$ for some $i$} \Return the leaf for the least such $i$ \EndIf
  \For{$k \gets 2$ \textbf{to} $n$}
    \For{$S$ with $|S| = k$ in increasing mask order}
      \If{$\textsc{Find}(S, t) \neq \bot$} \Return $(t, \hat\rho(\tau(S, t)))$ \EndIf
    \EndFor
    \If{$k < n$} \For{$S$ with $|S| = k$} $\textsc{Build}(S)$ \EndFor \EndIf
  \EndFor
  \State $(S_1, v_1) \gets$ least $\kappa$ over all stored values of subsets $S \ne [n]$
  \State $v_2 \gets$ least $(|t - v|, v)$ over results of the pairs of every split of $[n]$, streamed and not stored
  \State \Return $([n], v_2)$ if $|t - v_2| < |t - v_1|$, otherwise $(S_1, v_1)$, rendered through the origins
\EndProcedure
\end{algorithmic}
\end{algorithm}

\begin{theorem}[Correctness]\label{thm:correct}
Algorithm~\ref{alg:reference} returns $\operatorname{Solve}(t, p)$.
\end{theorem}

\begin{proof}
\emph{Build keeps the canonical derivation.} By induction on $|S|$ the stored lists are exactly the sets
$V(L)$, sorted. $\proc{Build}(S)$ then visits every derivation of every $v \in V(S)$ in increasing order of
the key $(\sigma_S(L), a, b, \operatorname{rank})$ of Definition~\ref{def:derivation}, and it keeps the first,
which is the least: this is $\omega(S, v)$. The set of keys is $V(S)$ by Theorem~\ref{thm:recursion}.

\emph{Find returns the canonical derivation.} The relation ``some legal operation on $a$ and $b$ yields $v$''
is symmetric in $a$ and $b$, so scanning either half finds the same set $H$, and Lemma~\ref{lem:partners}
lists, for each scanned value, every partner that could be in $H$. Splits are visited in order and the
least element of $H$ is taken, so the result is $\omega(S, v)$ whenever $v \in V(S)$, and $\bot$ otherwise.

\emph{The loop realises the objective.} Searching size $k$ needs only the lists for sizes below $k$, which
have been built. The first size and least mask with a hit give case~(a) of Proposition~\ref{prop:staged}. If
no size hits, every proper subset has been built, and the last three lines are case~(b) verbatim.

\emph{Rendering.} The stored origins are the canonical derivations, so following them from the answer
reproduces $\tau(S^*, v^*)$, which is rendered as $\hat\rho$.
\end{proof}

\paragraph{Cost.}
Table~\ref{tab:cost} measures the work by subset size. Building size $k$ forms every pair across every split
of every $k$-subset, while searching forms only one scan of the shorter half of each split. At size five,
searching scans 214 values and building forms 7{,}140 pairs, so the early exit of
Proposition~\ref{prop:early} matters far more than any constant factor. The search also replaces the
build for the whole pool, where building would form 15{,}456 pairs.

\begin{table}[ht]
\centering
\begin{tabular}{rrrr}
\toprule
$k$ & pairs to build size $k$ & legal results & values scanned to search size $k$ \\ \midrule
2 & 15 & 46 & 15 \\
3 & 184 & 571 & 60 \\
4 & 1{,}446 & 4{,}522 & 185 \\
5 & 7{,}140 & 22{,}334 & 214 \\
6 & 15{,}456 & 48{,}213 & 196 \\ \bottomrule
\end{tabular}
\caption{Work per round summed over all subsets of size $k$, averaged over the first 300 benchmark rounds.
Each scanned value costs at most six lookups.}\label{tab:cost}
\end{table}

\subsection{Refinements on the CPU}\label{sec:cpu}

The compiled implementations keep Algorithm~\ref{alg:reference} and change how $\proc{Build}$ removes
duplicates and sorts. Each refinement is exact, so the output is unchanged.

\paragraph{Duplicates by bitset.}
Instead of a map, a bitset $B$ over $\{0, \dots, M(S)\}$ records which values have been seen. Each legal
result is written to the next output slot unconditionally, and the output pointer advances by one minus
the old bit. The first derivation of each value is therefore the one kept, as before, with no branch to
mispredict.

\paragraph{Sorting by rank.}
Once every result is in, the bitset itself is the sorted set.

\begin{lemma}[Rank]\label{lem:rank}
Split $B$ into 64-bit words $B_0, B_1, \dots$, and let $c_j = \sum_{i < j} \operatorname{popcount}(B_i)$. A value
$v \in B$ belongs at index $c_{\lfloor v / 64 \rfloor} + \operatorname{popcount}\bigl(B_{\lfloor v / 64 \rfloor}
\mathbin{\&} (2^{v \bmod 64} - 1)\bigr)$ of the ascending list.
\end{lemma}

\begin{proof}
The expression counts the members of $B$ below $v$: whole words before $v$'s word, then the lower bits of its own.
\end{proof}

This sorts the list in time linear in the number of values and words, with no comparisons. Sets with
few values use insertion sort instead.

\paragraph{A tighter value bound.}
The bitset is sized by $M(S)$, where $M(\{i\}) = p_i$ and
$M(S) = \max_{\{L, R\}} \max\bigl(M(L) + M(R),\ M(L)\, M(R)\bigr)$. Every $v \in V(S)$ satisfies $v \le M(S)$, since
subtraction and division never exceed the larger operand and addition and multiplication are monotone.
This is at most the bound of Proposition~\ref{prop:bound}.

\paragraph{Division by reciprocal.}
Integer division is slow, so the solvers test exactness with a precomputed floating-point reciprocal.

\begin{lemma}[Reciprocal division]\label{lem:reciprocal}
Let $1 \le y \le x < 2^{31}$ and let $q$ be $x \cdot \operatorname{fl}(1/y)$, computed in IEEE double precision and
rounded to the nearest integer. Then $y \mid x$ if and only if $q y = x$.
\end{lemma}

\begin{proof}
With unit roundoff $u = 2^{-53}$, the computed product differs from $x / y$ by at most
$(x / y)(2u + u^2) < 2^{31} \cdot 2^{-51} < \tfrac12$. If $y \mid x$ then $x / y$ is the nearest integer, so
$q = x / y$ and $q y = x$. If $y \nmid x$, no integer $q$ satisfies $q y = x$.
\end{proof}

Values stay below $2^{27}$ (Proposition~\ref{prop:bound}), so the lemma applies with room to spare. An inexact
division writes $0$, whose bit is set before building starts, so the bitset drops it without a branch.

\paragraph{Known larger operand.}
Because $V(R)$ is sorted, the right values not exceeding a left value $a$ form a prefix. Looping over that
prefix and then the rest fixes which operand is larger without comparing them.

\subsection{Parallel algorithms on the GPU}\label{sec:gpu}

A GPU runs thousands of threads with no common order of events, so ``keep the first derivation found''
has no direct meaning. Two observations restore determinism.

\begin{proposition}[Derivation by minimum]\label{prop:min}
Encode a derivation $(L, a, b, \circ)$ of $v$ from $S$ as the integer
$\bigl(\sigma_S(L), a, b, \operatorname{rank}(\circ)\bigr)$ packed most significant first. Then
$\omega(S, v)$ is the minimum of the codes of all derivations, whatever order they are found in.
\end{proposition}

\begin{proof}
Packing with enough bits per field turns the lexicographic order of Definition~\ref{def:derivation} into
integer order, and the minimum of a set does not depend on the order of its elements.
\end{proof}

So threads can fold derivations into a result with an atomic minimum. The second observation is that
origins need not be stored at all.

\begin{proposition}[Sets suffice]\label{prop:sets}
Given the sets $V(L)$ for all $L \subsetneq S$, $\proc{Find}(S, v)$ computes $\omega(S, v)$ for any $v \in V(S)$.
Hence $\tau(S^*, v^*)$ can be recovered from the value sets alone with at most $k^* - 1$ calls to
$\proc{Find}$, one per internal node.
\end{proposition}

\begin{proof}
The first claim is the second step of the proof of Theorem~\ref{thm:correct}, which used only the sets. The tree
has $k^*$ leaves and so $k^* - 1$ internal nodes, each of which is the canonical derivation of its own subset
and value.
\end{proof}

The CUDA solver uses both. One thread block solves one round in on-chip shared memory. It stores each
subset's values in a hash table keyed by subset and value, so duplicates are removed by an atomic
compare-and-swap and nothing is sorted. The nodes of the answer tree at the same depth are independent,
so they are expanded one level at a time, with all threads searching at once. The PyTorch solver instead
keeps every value of every round in one sorted array of keys (size, round, mask, value) and answers
membership queries by binary search. It removes duplicates with a stable sort of results generated in
canonical order, so the first derivation of each value survives.

Both GPU solvers also avoid building the subsets of size $n - 1$ when the answer uses every number.

\begin{lemma}[Deferred level]\label{lem:defer}
Let $n \ge 3$. In any split of $[n]$, if one half has size $n - 1$ then the other is a singleton. So
$\proc{Find}([n], t)$ can run with the sets of size $n - 1$ unbuilt: in those splits it scans the singleton
and decides each of at most six partners $b$ by a nested $\proc{Find}(L, b)$ with $|L| = n - 1$, which needs
only sets of size at most $n - 2$.
\end{lemma}

\begin{proof}
The sizes of the two halves add up to $n$. Nested searches return membership correctly by the second step of
Theorem~\ref{thm:correct}.
\end{proof}

There are $n$ such splits with at most six partners each, so at most $6n = 36$ nested searches, against
7{,}140 pairs to build size five. The sets of size $n - 1$ are still built when the whole pool misses the
target, because case~(b) of Proposition~\ref{prop:staged} needs every proper subset. By the same counting,
a subset of size $n - 2$ is a half only in splits of sets of size $n - 1$ or $n$, where the other half has size
$1$ or $2$. When $n \ge 5$ that other half is the smaller one, so searches scan it and only look up the
values of the $(n - 2)$-subset. The CUDA solver therefore keeps those sets as hash tables without a list,
listing them by a table scan only in the rare rounds that build size $n - 1$. This cuts its shared memory
enough to run three rounds on each multiprocessor at once. For $n = 4$ the argument fails, since the split
$2 + 2$ has equal halves, and the solver lists those sets as usual.

\subsection{Precomputed and learned answers}\label{sec:approx}

Two further versions trade exact text for speed. Their answers are equally good but not always identical.

\paragraph{Lookup table.}
By Proposition~\ref{prop:pools} the game has only $8{,}548{,}200$ instances up to the order of the pool, so
their answers can all be computed in advance. The table is keyed by the sorted pool, so by
Remark~\ref{rem:invariance} it stores $\operatorname{Solve}$ of a permutation of the dealt pool: the same
$(d^*, k^*)$, possibly with different text. Keying on dealt order instead would need about 430 million
answers.

\paragraph{Learned proposals with a certificate.}
A small transformer reads $(t, p)$ and proposes an expression one token at a time, with decoding constrained so
that every proposal is legal. A proposal is only a hint, and the solver accepts it on the following
certificate.

\begin{proposition}[Certificate]\label{prop:certificate}
Suppose a legal expression $e$ over $p$ has value $t$ and uses $k$ numbers, and $\proc{Find}(S, t) = \bot$ for
every $S$ with $|S| < k$. Then $d^* = 0$, $k^* = k$ and $e$ is equally good as $\operatorname{Solve}(t, p)$.
\end{proposition}

\begin{proof}
The expression shows that $t \in V(\leaves(e))$, so $d^* = 0$ and $k^* \le k$, and the failed searches show
$k^* \ge k$.
\end{proof}

Checking the certificate needs only subsets of size up to $k - 2$ built, so an accepted proposal saves
building size $k - 1$. A rejected or missing proposal falls back to the full algorithm, so the hybrid is
never worse than the solver. In practice the model proposes a shortest answer too rarely, on 7\% of
six-number rounds, for the saving to outweigh the cost of running it.

\section{Empirical facts}\label{sec:data}

The benchmark set \texttt{bench/data/rounds10k.txt} holds $10{,}000$ game-distributed rounds. Its
reference answers \texttt{ref10k.tsv} are $\operatorname{Solve}(t, p)$ as computed by the Rust solver.

\begin{table}[h]
\centering
\begin{tabular}{lr}
\toprule
Outcome & Rounds \\ \midrule
Exact ($d^* = 0$) & 9{,}776 \\
\quad using 1 number & 3 \\
\quad using 2 numbers & 81 \\
\quad using 3 numbers & 860 \\
\quad using 4 numbers & 3{,}622 \\
\quad using 5 numbers & 4{,}063 \\
\quad using 6 numbers & 1{,}147 \\
Closest at $d^* = 1$ & 207 \\
Closest at $d^* = 2$ & 13 \\
Closest at $d^* = 3$ & 2 \\
Closest at $d^* = 4$ & 2 \\ \bottomrule
\end{tabular}
\caption{Optimal distance and count over the 10{,}000 benchmark rounds.}
\end{table}

\begin{table}[h]
\centering
\begin{tabular}{rrr}
\toprule
$|S|$ & mean $|V(S)|$ & largest $|V(S)|$ \\ \midrule
1 & 1.0 & 1 \\
2 & 3.1 & 4 \\
3 & 17.1 & 28 \\
4 & 133.8 & 228 \\
5 & 1{,}079.1 & 2{,}181 \\
6 & 8{,}154.1 & 18{,}946 \\ \bottomrule
\end{tabular}
\caption{Sizes of reachable sets over the first 300 benchmark rounds. The largest value seen was 25{,}200{,}000.}\label{tab:sets}
\end{table}

Two consequences shape the solvers. First, almost every round is exact, so by
Proposition~\ref{prop:early} the full pool's value set is almost never needed. The five-number sets are needed only by the 11\% of
rounds whose answer uses all six numbers, and the GPU solvers avoid even those by Lemma~\ref{lem:defer},
building them only for the 2\% of rounds that miss. Second, the work is dominated by the largest subsets. The reachable sets
grow roughly eightfold with each extra number, and most of the cost lies in combining two halves of a
five- or six-number subset.

\section{Performance}\label{sec:performance}

Every implementation was timed on the same 10{,}000 benchmark rounds. Table~\ref{tab:perf} lists the
results, copied from the repository's \texttt{PERFORMANCE.md}, fastest first.

\paragraph{Method.}
The machine is an Intel Core i9-13900HX with an NVIDIA RTX 4070 Laptop GPU, running Linux 7.0. CPU
benchmarks are pinned to one performance core. \emph{Solve time} is the best of three runs, measured
inside a harness for each language, so that start-up and output are excluded. Python and bwBASIC are timed on a
1{,}000-round set because the full set takes over a minute. The just-in-time compiled languages are timed after three
untimed warm-up passes. The GPU harnesses warm up first too, so that CUDA start-up, kernel compilation and
the GPU's climb from its idle clock are excluded. \emph{Straight port} is the solve time of each language's
first faithful translation, before any language-specific optimisation. \emph{Whole program} is the mean
wall-clock time of each language's own executable dealing and solving one random round, over 300 runs (10
for the two PyTorch versions), so it includes process and runtime start-up.

The twelve newest languages (FreeBASIC, BASIC, WebAssembly, F\#, Haskell, Kotlin, Common Lisp, OCaml, Ruby,
Julia, Erlang and Elixir) were measured later on 2026-10-09, in sessions in which C measured 0.0182--0.0185\,ms, and ISPC on 2026-10-10, when C measured 0.0183\,ms. The four GPU and precomputed rows were measured a day after the original set, in a quieter session in which C
measured 0.0184--0.0188\,ms rather than 0.0208\,ms. Comparisons between those rows and the others
therefore carry roughly 10\% uncertainty. The GPU rows are throughput, with all 10{,}000 rounds submitted as
one batch, not the latency of a single round.

\begin{table}[p]
\centering
\footnotesize
\begin{tabular}{@{}lrrrrl@{}}
\toprule
Implementation & \multicolumn{1}{c}{Solve} & \multicolumn{1}{c}{Straight port} & \multicolumn{1}{c}{Gain}
  & \multicolumn{1}{c}{Whole program} & Answers \\
 & ms/round & ms/round & & ms/run & \\ \midrule
Rust + CUDA (GPU, batched)          & 0.000466 & 0.0107\rlap{$^a$} & 23$\times$  & 203   & identical \\
Lookup table (Python)               & 0.00069  & --                & --          & 13.4  & 31\% ties \\
PyTorch on the GPU (batched)        & 0.0157   & 0.142\rlap{$^a$}  & 9.0$\times$ & 3{,}218 & identical \\
ISPC (AVX2, four 64-bit lanes)\rlap{$^d$} & 0.0115 & 0.078       & 6.8$\times$ & 0.35  & identical \\
Zig (ReleaseFast)                   & 0.0203   & 0.083             & 4.1$\times$ & 0.48  & identical \\
C (Clang)                           & 0.0208   & 0.118             & 5.7$\times$ & 0.39  & identical \\
Julia 1.13                          & 0.0208   & 0.259             & 12.5$\times$ & 548   & identical \\
C++ (Clang)                         & 0.0213   & 0.086             & 4.0$\times$ & 0.73  & identical \\
FreeBASIC (GCC back end)            & 0.0215   & 0.123             & 5.7$\times$ & 0.50  & identical \\
x86-64 assembly                     & 0.0218   & 0.043\rlap{$^a$}  & 2.0$\times$ & 0.23  & identical \\
WebAssembly (C in Node)             & 0.0229   & 0.243             & 10.6$\times$ & 17.4 & identical \\
Fortran (gfortran)                  & 0.0247   & 0.109             & 4.4$\times$ & 0.52  & identical \\
Ada (GNAT)                          & 0.0260   & 0.260             & 10$\times$  & 0.78  & identical \\
C\# (.NET 10)                       & 0.0274   & 0.212             & 7.7$\times$ & 21.9  & identical \\
Swift                               & 0.0274   & 0.212             & 7.7$\times$ & 1.52  & identical \\
F\# (.NET 10)                       & 0.0287   & 0.332             & 11.6$\times$ & 30.3 & identical \\
Go                                  & 0.0330   & 0.190             & 5.8$\times$ & 0.64  & identical \\
Pascal (Free Pascal)                & 0.0352   & 0.752             & 21$\times$  & 0.38  & identical \\
PyTorch with transformer proposals  & 0.0354   & --                & --          & 3{,}223 & 8\% ties \\
Haskell (GHC)                       & 0.0360   & 0.639             & 17.8$\times$ & 0.82  & identical \\
Kotlin (JVM)                        & 0.0373   & 0.149             & 4.0$\times$ & 48.2  & identical \\
Java (OpenJDK 25)                   & 0.0391   & 0.154\rlap{$^b$}  & 3.9$\times$ & 34.2  & identical \\
Common Lisp (SBCL)                  & 0.0472   & 0.446             & 9.4$\times$ & 1.51  & identical \\
OCaml                               & 0.0537   & 0.450             & 8.4$\times$ & 0.64  & identical \\
Rust                                & 0.0557   & 9.7\rlap{$^a$}    & 174$\times$ & 0.51  & identical \\
Erlang (BEAM)\rlap{$^c$}            & 0.0575   & 0.494             & 8.6$\times$ & 73.7  & identical \\
Elixir (BEAM)\rlap{$^c$}            & 0.0576   & 0.729             & 12.7$\times$ & 83.6  & identical \\
JavaScript (Node 24)                & 0.0848   & 0.381             & 4.5$\times$ & 24.3  & identical \\
COBOL (GnuCOBOL)                    & 0.173    & 2.07              & 12$\times$  & 1.30  & identical \\
LuaJIT                              & 0.186    & 0.61              & 3.3$\times$ & 2.11  & identical \\
Ruby with RJIT                      & 0.291    & --                & --          & 29.0  & identical \\
Forth (\texttt{gforth-fast})        & 0.318    & 0.864             & 2.7$\times$ & 4.28  & identical \\
Forth (\texttt{gforth})             & 0.399    & 1.036             & 2.6$\times$ & 7.02  & identical \\
Ruby 3.3                            & 0.494    & 2.22              & 4.5$\times$ & 29.6  & identical \\
Lua 5.4                             & 0.550    & 1.10              & 2.0$\times$ & 2.44  & identical \\
Python                              & 7.23     & 12.0\rlap{$^a$}   & 1.7$\times$ & 17.6  & 1.5\% ties \\
BASIC (bwBASIC, line-numbered)      & 306      & 1{,}114        & 3.6$\times$ & 282   & identical \\ \bottomrule
\end{tabular}
\caption{Solve time per round over the 10{,}000 benchmark rounds (1{,}000 for Python and bwBASIC), the time before optimisation, and the time
for a whole program run. $^a$First version rather than a straight port of another language. $^b$Re-measured with the current harness (Section~\ref{sec:surprises}); 0.257\,ms was recorded originally. $^c$Defers the four- and five-number subsets, extending Lemma~\ref{lem:defer}, and so does less work than the other CPU versions; Erlang built like the others takes about 0.22\,ms. $^d$Measured on 2026-10-10 when C took 0.0183\,ms, so 1.6 times faster than C (Section~\ref{sec:surprises}). ``Ties'' are
equally good answers that differ in text from $\operatorname{Solve}$; no implementation gives a worse answer.}
\label{tab:perf}
\end{table}

\paragraph{Observations.}
\begin{enumerate}[label=(\arabic*)]
\item \emph{The algorithm matters more than the language.} The compiled languages from Zig to Java lie within a
factor of two of one another, while every one of them gained between two and twenty-one times from the
techniques of Sections~\ref{sec:reference} and~\ref{sec:cpu}. The largest single gain, Rust's
174$\times$, came from moving its first version to the staged search with early exit
(Proposition~\ref{prop:early}) and searching rather than building (Lemma~\ref{lem:partners}), together with a compact bitset store.
\item \emph{Compilers set the floor.} C, C++, Zig, FreeBASIC and hand-written assembly finish within 8\% of one another,
because they share one design and differ only in code generation. Garbage-collected and managed languages
come close when written over flat buffers, and interpreters trail C by between nine times and, for the line-numbered bwBASIC, about fifteen thousand times. WebAssembly in Node comes within 25\% of native C. The one solver below that floor, ISPC, changes the design to suit SIMD lanes (Section~\ref{sec:surprises}).
\item \emph{Parallelism needs batches.} The CUDA solver is about 40 times faster than C per round, but only
when thousands of rounds are in flight, one per thread block. The PyTorch solver takes 11\,ms for a single
round, slower than plain Python, and only passes C at around 10{,}000 rounds per batch. Their whole-program
times, 0.2 and 3.2 seconds, are almost entirely runtime and GPU start-up. The parallel formulation of
Section~\ref{sec:gpu} makes the GPU versions exact, but it does not make them suitable for one round at a time. Sharing one round across CPU threads gains at most a factor of two (Section~\ref{sec:parallel}).
\item \emph{Precomputation wins on latency.} Because the instance space is finite
(Proposition~\ref{prop:pools}), a 248\,MB table answers in 0.7\,\textmu s from plain Python, faster than any
solver except CUDA. The price is that its text matches $\operatorname{Solve}$ for only 69\% of rounds
(Remark~\ref{rem:invariance}).
\item \emph{Learning did not pay.} The transformer's certified proposals (Proposition~\ref{prop:certificate})
never give a wrong answer, but they too rarely save the expensive levels to cover the cost of running the
model, so the hybrid is 2.3 times slower than the PyTorch solver it guides.
\end{enumerate}

\section{Straight ports against optimised versions}\label{sec:optimisation}

Every language was implemented twice. The \emph{straight port} is a faithful translation of
Algorithm~\ref{alg:reference} as the Rust reference writes it. It uses the language's ordinary data
structures, the library sort, a branching binary search and the \texttt{\%} operator for the division test.
The \emph{optimised version} keeps the same algorithm and the same output, and changes only how the work is
done. Both versions of every language were checked against the reference answers. Table~\ref{tab:opt}
summarises what changed, and the rest of this section explains why the changes help.

\subsection{Changes shared across languages}

Most optimised versions adopt some or all of the refinements of Section~\ref{sec:cpu}. In rough order of
impact they are:
\begin{enumerate}[label=(\alph*)]
\item a 16-byte entry holding the value, operator, masks and both operands, written without a branch and
      kept only if the bitset shows the value is new;
\item the reciprocal division test of Lemma~\ref{lem:reciprocal} in place of an integer division;
\item each pair loop split at the left value, so the larger operand is known without a comparison;
\item sorting by rank in the bitset (Lemma~\ref{lem:rank}), which needs no comparisons;
\item branchless binary search and other unpredictable comparisons turned into conditional moves or
      arithmetic masks;
\item the bitset sized by the bound $M(S)$, so it grows only when larger subsets are built.
\end{enumerate}
Each of these is an exact transformation of $\proc{Build}$ or $\proc{Find}$, so none can change
$\operatorname{Solve}$. The straight ports already had the staged search, the early exit and the
partner lookups, because the Rust reference they copy has them.

\begin{table}[p]
\centering
\scriptsize
\setlength{\tabcolsep}{4pt}
\begin{tabular}{@{}lrrrp{9.5cm}@{}}
\toprule
Language & \multicolumn{1}{c}{Straight} & \multicolumn{1}{c}{Optimised} & \multicolumn{1}{c}{Gain} & Principal changes \\
 & ms/round & ms/round & & \\ \midrule
ISPC & 0.078 & 0.0115 & 6.8$\times$ & values only, deduplicated by bit OR; pair loops across four 64-bit lanes; lists read from bitsets; bit-test partner checks \\
Zig & 0.083 & 0.0203 & 4.1$\times$ & C allocator instead of the default, which returned memory to the kernel on every free; buffers grown by doubling; branch hint for a conditional move in the search \\
C & 0.118 & 0.0208 & 5.7$\times$ & all of (a)--(f); Clang instead of GCC; binary search adds \texttt{(value <= target) * half}; subsets listed in size order up front \\
Julia & 0.259 & 0.0208 & 12.5$\times$ & values listed straight from the bitset with \texttt{trailing\_zeros}, so no sort; an \texttt{llvmcall} select marked \texttt{!unpredictable} to force conditional moves \\
C++ & 0.086 & 0.0213 & 4.0$\times$ & the C design with templates and lambdas; buffers that skip zero-filling; Clang instead of GCC \\
FreeBASIC & 0.123 & 0.0215 & 5.7$\times$ & port of the optimised C with typed pointers and macros; \texttt{-march=native} for a one-instruction popcount; \texttt{IIf} replaced by a mask \\
Assembly & 0.043$^a$ & 0.0218 & 2.0$\times$ & rank sort; branchless stores; conditional moves in the search; branchless division test \\
WebAssembly & 0.243 & 0.0229 & 10.6$\times$ & the optimised C compiled unchanged with portable stand-ins for three intrinsics; quotient rounded by adding $1.5 \cdot 2^{52}$ \\
Fortran & 0.109 & 0.0247 & 4.4$\times$ & \texttt{-O3} so the store is inlined; no default component values, which zeroed whole arrays; arrays passed separately so they cannot alias; masks instead of \texttt{merge} \\
Ada & 0.260 & 0.0260 & 10$\times$ & runtime checks suppressed with \texttt{-gnatp}; cross-unit inlining; popcount and rounding intrinsics; modular type for masks \\
C\# & 0.212 & 0.0274 & 7.7$\times$ & near line-for-line port of the optimised C using \texttt{unsafe} pointers, native memory and \texttt{PopCount} \\
Swift & 0.212 & 0.0274 & 7.7$\times$ & unsafe pointer buffers; wrapping arithmetic; \texttt{-Ounchecked}; ternaries replaced by arithmetic masks \\
F\# & 0.332 & 0.0287 & 11.6$\times$ & port of the optimised C\# over native memory; conversions through signed \texttt{int64}; arithmetic mask in the division test \\
Go & 0.190 & 0.0330 & 5.8$\times$ & flat harness data so the garbage collector stops rescanning; \texttt{GOAMD64=v3}; masks instead of branches; lookup trimmed to fit the inliner \\
Pascal & 0.752 & 0.0352 & 21$\times$ & C allocator via \texttt{cmem}; 64-bit floating constants so the reciprocal is exact \\
Java & 0.154$^b$ & 0.0391 & 3.9$\times$ & entries packed into a \texttt{long[]}; \texttt{Long.bitCount} rank sort; straight-line partner checks \\
Kotlin & 0.149 & 0.0373 & 4.0$\times$ & port of the optimised Java with \texttt{LongArray} entries \\
Haskell & 0.639 & 0.0360 & 17.8$\times$ & \texttt{ST} over unboxed \texttt{PrimArray}s; bitset with rank sort; values only, origins rebuilt for the answer \\
Common Lisp & 0.446 & 0.0472 & 9.4$\times$ & port of the optimised C with typed arrays, \texttt{logcount} and \texttt{(safety 0)} \\
OCaml & 0.450 & 0.0537 & 8.4$\times$ & byte bitset with radix sort; integer-only \texttt{max} and \texttt{min}; \texttt{Bigarray}s to spare the collector; values only \\
Rust & 9.7$^a$ & 0.0557 & 174$\times$ & staged search with early exit; values only while searching; compact entries with a bitset; partner lookups; none of (a)--(f) yet \\
Erlang$^c$ & 0.494 & 0.0575 & 8.6$\times$ & rewritten to Elixir's design, with tuples for binary search and subset tables as constants \\
Elixir$^c$ & 0.729 & 0.0576 & 12.7$\times$ & values only with \texttt{:lists.usort}; four- and five-number subsets deferred; ordered partner checks; outward search for closest rounds \\
JavaScript & 0.381 & 0.0848 & 4.5$\times$ & four 32-bit words per entry in a \texttt{Uint32Array}; 32-bit bitset words with a hand-written popcount; no typed-array views in the insertion sort \\
COBOL & 2.07 & 0.173 & 12$\times$ & every hot operation kept out of GMP decimal arithmetic: shift-and-add multiplication, long division, a byte stamp per value and a byte radix sort \\
LuaJIT & 0.61 & 0.186 & 3.3$\times$ & each origin packed into one number; a radix sort written in Lua; hash tables pre-sized \\
Ruby & 2.22 & 0.494 & 4.5$\times$ & a Hash per subset with \texttt{key?} checks; split pair loops; whole-pool pass unrolled \\
Forth (fast) & 0.864 & 0.318 & 2.7$\times$ & byte stamp per value; byte radix sort; pointer stepping; hot words compiled in place \\
Forth & 1.036 & 0.399 & 2.6$\times$ & as above \\
Lua 5.4 & 1.10 & 0.550 & 2.0$\times$ & each origin packed into one number; \texttt{table.sort} kept \\
Python & 12.0$^a$ & 7.23 & 1.7$\times$ & values without origins while searching; origins rebuilt only for the answer \\
bwBASIC & 1{,}114 & 306 & 3.6$\times$ & one hash table of values without origins; hot subroutines placed first because \texttt{GOTO} searches from the top \\ \midrule
PyTorch & 0.142$^a$ & 0.0157 & 9.0$\times$ & expressions typeset on the GPU; pair generation fused by \texttt{torch.compile}; shorter half scanned; five-number subsets deferred; illegal results dropped before sorting \\
Rust + CUDA & 0.0107$^a$ & 0.000466 & 23$\times$ & split tables computed once; each phase flattened across all threads; four-number subsets kept as tables only; values appended as they are found; answer unfolded a level at a time \\ \bottomrule
\end{tabular}
\caption{Straight port against optimised version for each language. $^a$First version rather than a
straight port. $^b$Re-measured (Section~\ref{sec:surprises}). $^c$Defers the four- and five-number subsets, extending Lemma~\ref{lem:defer}, and so does less work than the other CPU versions.}\label{tab:opt}
\end{table}

\subsection{What the gains reveal}

The changes fall into a handful of recurring kinds, and which kinds matter says more about each language's
implementation than about the language itself.

\paragraph{Memory management.}
Several of the largest single gains had nothing to do with arithmetic. Zig's default allocator and Free
Pascal's heap manager both returned large blocks to the kernel on every free. Since each solve allocates
its own memory, Zig's version took 255{,}000 page faults over a run. Switching to the C library's allocator halved
both run times. Go spent a third of its time in the garbage collector rescanning the harness's 20{,}000
small objects, and gfortran zeroed whole arrays because of default component values. In each case the
straight port was paying for the runtime rather than the algorithm.

\paragraph{Branches.}
The binary search in $\proc{Find}$ and the duplicate test in $\proc{Build}$ take each direction about
half the time, so a branch there mispredicts constantly. Every compiled language gained by making them
branchless, but each needed a different way to ask for it. Zig has an explicit hint. Clang ignored its
hint, so the C version multiplies by a comparison. Go has no way to ask, so its version uses arithmetic
masks, and so do Fortran and Swift, whose compilers turned \texttt{merge} and the ternary operator back
into branches.

\paragraph{Slow primitives.}
Integer division is slow, so the reciprocal test of Lemma~\ref{lem:reciprocal} pays off in every
compiled language. Free Pascal first stored the constant $1.0$ in single precision, which made the
reciprocal too inexact for the lemma to hold until the precision was raised. COBOL is the extreme case.
GnuCOBOL compiles multiplication, division and bitwise operations into calls to GMP's decimal arithmetic,
which took 65\% of the time. The optimised version restricts its hot paths to the operations that compile to
native code and builds the rest from them, which cut decimal arithmetic to 3\%.

\paragraph{Sorting.}
Library sorts compare, and comparisons through a function pointer are slow. This is why straight C, which uses
\texttt{qsort}, started slower than straight C++, whose \texttt{std::sort} is inlined for the entry type. The
rank sort of Lemma~\ref{lem:rank} removes comparisons entirely where a population count is cheap. COBOL and Forth
use a byte radix sort instead, and so does LuaJIT, which cannot compile \texttt{table.sort}. Lua 5.4 kept \texttt{table.sort}, because its
interpreter runs a radix sort written in Lua three times slower than the built-in.

\paragraph{Safety checks.}
Ada's straight port keeps the language's default bounds and overflow checks, and its optimised build
suppresses them alongside inlining and native code generation. Swift's
\texttt{-Ounchecked} took it from 0.043 to 0.030\,ms. Removing checks is only safe once the code is known to be
in range, so the optimised Ada solver was also run over all 10{,}000 rounds with checks enabled, with
identical answers.

\paragraph{Interpreters.}
In Lua and Forth the cost tracks the number of operations the interpreter dispatches, not the machine work
behind them. The gains there come from doing less per value: packing an origin into one number instead of a
table, stamping a byte with the subset's mask instead of clearing a bitset, and stepping pointers rather than
computing indices.

\paragraph{Line-numbered BASIC.}
bwBASIC interprets the program text almost directly. It costs about 2.4\,\textmu s per statement, plus about
1\,\textmu s per operator, and a fifth of its time is spent in \texttt{strcmp}. It finds every \texttt{GOTO} and
\texttt{GOSUB} target by searching from the start of the program. A jump to a line near the end therefore costs
about 3\,\textmu s, against 0.4\,\textmu s near the start, and the optimised solver puts its hot subroutines
first, so that line numbers themselves become a performance decision. Its 3.6-fold gain otherwise came from
Proposition~\ref{prop:sets}: it keeps one hash table of values, with no sorting and no stored origins, and
rebuilds the few origins the answer needs by search.

\paragraph{Python.}
The Python version took a different route, using Proposition~\ref{prop:sets}. It stores only the value sets while
searching and rebuilds origins only for the nodes of the final answer. That saves about 40\%. The interpreter's
cost per value still dominates, so Python remains 350 times slower than C.

\paragraph{What did not carry over.}
A SIMD rewrite of the assembly was ruled out because removing duplicates is serial: whether a value is new
depends on every value before it. Profile-guided optimisation gave at most 3\% in C, Go and Fortran. Radix sort
lost to rank sort in C at every sparsity tested. The GPU versions are a different story. Their gains come
from restructuring for parallelism (Section~\ref{sec:gpu}) rather than from any of (a)--(f).

\subsection{Surprises examined}\label{sec:surprises}

Adding Haskell, OCaml, Common Lisp, F\#, Kotlin, Ruby, WebAssembly, Erlang, Elixir and ISPC produced several results that go
against the usual reputations of these languages. This subsection gives the evidence for each one and our
reading of the cause. Times are best of three, pinned to one core. Hardware counters come from
\texttt{perf stat}. To leave out start-up and input parsing, they difference a run of six passes over the
10{,}000 rounds against a run of one pass. Where we infer a cause rather than measure it, we say so.

\begin{table}[ht]
\centering
\small
\begin{tabular}{@{}lrrrrrr@{}}
\toprule
Language & ms/round & instructions & cycles & IPC & branch misses & loads \\ \midrule
C (Clang)          & 0.0187 & 324{,}000 & 93{,}400  & 3.47 & 797   & 52{,}300 \\
Go                 & 0.0303 & 541{,}000 & 154{,}700 & 3.50 & 985   & 106{,}300 \\
Haskell (GHC)      & 0.0366 & 596{,}000 & 186{,}500 & 3.20 & 1{,}576 & 140{,}800 \\
Common Lisp (SBCL) & 0.0487 & 800{,}000 & 237{,}500 & 3.37 & 845   & 168{,}800 \\
Rust               & 0.0503 & 608{,}000 & 259{,}900 & 2.34 & 3{,}665 & 134{,}400 \\
OCaml              & 0.0541 & 1{,}102{,}000 & 272{,}500 & 4.04 & 1{,}676 & 203{,}400 \\ \bottomrule
\end{tabular}
\caption{Hardware counters per round for the ahead-of-time compiled solvers, on the 10{,}000 benchmark
rounds.}\label{tab:counters}
\end{table}

\paragraph{Common Lisp beats Rust.}
SBCL is a dynamically typed Lisp, yet its optimised solver, at 0.049\,ms, beats the Rust solver, at 0.050\,ms. The
counters show why. Rust executes \emph{fewer} instructions than Lisp, but it suffers 3{,}665 branch
mispredictions per round against Lisp's 845, which is close to C's 797. At roughly 15--20 cycles each, the
extra 2{,}800 mispredictions cost about 50{,}000 cycles. That is most of the difference, and it drags Rust's
instructions per cycle (IPC) down to 2.34. The Lisp version is a direct port of the optimised C design, including the
branchless append and binary search of refinements (a) and (e). The Rust solver has not yet had any of
(a)--(f). So this result is about design rather than language: SBCL with full type declarations and
\texttt{(safety 0)} compiles to plain machine arithmetic, and the branchless design matters more than the
better compiler. Profiling also shows that 11--15\% of the Lisp solver's cycles are in \texttt{memset}, called from
SBCL's allocator (\texttt{lisp\_alloc}). The allocator zeroes each fresh region of memory it hands out,
and the solver allocates its arrays anew for every solve. Lisp therefore has further headroom.

\paragraph{Haskell comes within a factor of two of C.}
A lazy, purely functional language might be expected to trail far behind, but the optimised Haskell solver
runs at 0.037\,ms. That puts it between Go and Java, and ahead of OCaml, Lisp and Rust. Two-thirds of its cycles
are in a single function, \texttt{\$wsolve}. This is GHC's \emph{worker}: after strictness analysis
has shown that every argument will be evaluated, GHC splits off a version that takes them unboxed, in
registers. The state lives in unboxed \texttt{PrimArray}s inside the \texttt{ST} monad, which compile to
plain loads and stores, and the garbage collector's share is under 5\%. What remains is a code-generation
gap. Haskell executes 1.84 times C's instructions and suffers twice C's branch mispredictions. We think,
without having checked the assembly, that GHC's native code generator emits real branches where Clang
emits conditional moves.

\paragraph{OCaml trails Haskell.}
OCaml is usually the faster of the two, but here it is the slowest of the compiled functional languages.
Its counters are unusual. It runs at 4.04 instructions per cycle, the best of any solver, so the processor
is not stalling. It simply executes 3.4 times as many instructions as C. We attribute this to three
causes. First, OCaml integers are 63 bits with a tag bit, so most arithmetic needs extra
instructions to adjust the tag. Second, OCaml has no population-count primitive, so the rank sort of
Lemma~\ref{lem:rank} is unavailable: the solver keeps a byte per 8 values and sorts with a multi-pass radix sort
instead. Third, the compiler installed here lacks the \emph{flambda} optimiser, so it inlines less.
Garbage collection costs about 3\% after the large arrays were moved to \texttt{Bigarray}s. Haskell wins because GHC
can work on raw unboxed integers, which suits an integer-array workload like this one.

\paragraph{WebAssembly is 3.3 times faster than JavaScript in the same engine.}
Both run in V8, yet the C solver compiled to WebAssembly takes 0.023\,ms, against 0.076\,ms for the
hand-optimised JavaScript solver. That is only 1.24 times native C. The difference is in what each format
lets the engine assume. JavaScript's bitwise operators work on 32 bits, so the JavaScript bitset has twice
as many words and counts bits with shifts and masks. Every number must also be checked for its type at run
time. WebAssembly has native 64-bit integers and a \texttt{popcnt} instruction, its types are fixed in
advance, and Clang has already optimised the code before V8 sees it. On 64-bit hosts V8 also
avoids most bounds checks on WebAssembly memory accesses by reserving guard pages around them. The
remaining 24\% over native comes from about 1.5 times native's instruction count. Neither BMI2 nor the
direct rounding conversion used by the C version is available, and V8's compiler optimises less
aggressively than Clang. The straight port suffers far more, running 2.1 times slower than its native
build. About half of its time goes to the C library's \texttt{qsort}, which calls the comparison function
through an indirect call. Another fifth goes to V8's out-of-line handling of the variable-length
\texttt{memory.copy} instructions behind \texttt{realloc}. The optimised solver uses neither sorting by
comparison nor frequent reallocation, which is why WebAssembly costs it so little.

\paragraph{F\# needed one change C\# did not.}
F\# and C\# share the .NET runtime and its compiler, and the optimised F\# solver, a port of the C\# design,
ends within 10\% of it. The largest single gain, from 0.032 to 0.0275\,ms, came from converting between
integers and floating point through \emph{signed} 64-bit integers. Before AVX-512, x86-64 has no instruction
that converts an unsigned 64-bit integer to or from a double, so the compiler emits a branching sequence of
several instructions for each conversion. The values stay below $2^{27}$
(Proposition~\ref{prop:bound}), so the signed conversion is exact. The straight ports differ more:
F\#'s takes 0.33\,ms and C\#'s 0.19\,ms. A profile of the F\# port puts 15\% of its cycles in runtime type
tests (\texttt{IsInstance}) and in F\#'s generic structural comparison. Idiomatic F\# compares tuples
structurally, for instance in \texttt{Seq.minBy} with a tuple key, and that comparison goes through a
generic comparer that discovers types at run time. F\#'s \texttt{option} type also allocates an object for
every \texttt{Some}, where C\# uses a null reference.

\paragraph{A surprise that was not.}
At first the Kotlin straight port, at 0.15\,ms, appeared much faster than Java's recorded 0.257\,ms. We
suspected Kotlin's \texttt{inline} functions, which remove the per-pair lambda calls Java's port makes.
Re-timing Java's straight port with the current harness gives 0.154\,ms, the same as Kotlin. The old figure
dates from before the harness ran three full warm-up passes, so it measured HotSpot part-way through compiling.
Table~\ref{tab:perf} and Table~\ref{tab:opt} now show the new figure. The optimised Kotlin and Java solvers also tie,
at 0.0374 and 0.0376\,ms, as expected of two languages compiling to near-identical bytecode.

\paragraph{Interpreted Ruby beats interpreted Lua.}
The optimised Ruby solver runs at 0.51\,ms in the standard interpreter, slightly ahead of Lua 5.4's 0.55\,ms.
It runs at 0.31\,ms under Ruby's RJIT compiler. Lua is normally considered the lighter and faster interpreter. The
profiles point to sorting. In Lua 5.4, 21\% of all cycles are in \texttt{lua\_geti}, \texttt{lua\_seti} and
\texttt{lua\_compare}, the general-purpose C interface that \texttt{table.sort} uses to read, write and
compare every element. Ruby's \texttt{Array\#sort} on integers instead calls the C library's merge sort with
a fast path for comparing small integers, which takes 4\% of its cycles. Both interpreters spend most of
their remaining time in hash-table lookups, where they are evenly matched.

\paragraph{Two languages on one virtual machine.}
Erlang and Elixir compile to the same BEAM bytecode and run on the same just-in-time compiler, yet their first
optimised solvers differed by a factor of 2.8: 0.162\,ms for Erlang against 0.058\,ms for Elixir. The
difference was entirely in design. Building a value costs 50--80\,ns on BEAM, against a few nanoseconds in C,
so it pays to build less even at the price of more searches. Elixir left both the four- and five-number
subsets unbuilt, extending Lemma~\ref{lem:defer} by one level, stopped at the first matching partner, and
searched outward from the target ($t \pm 1$, $t \pm 2$, up to 8) before building the whole pool in the 2\% of
rounds that miss. Those rounds had taken 38\% of its time. Rewritten to the same design, Erlang runs at
0.0575\,ms, level with Elixir. Removing each idea from the rewritten version in turn shows that the extra
deferred level gives the largest gain and the outward search about a quarter. Building none of the
subsets was the opposite of what PyTorch found, where deferring the four-number subsets as well cost 30\%.

\paragraph{SIMD lanes break the C floor.}
ISPC compiles a C-like program to run on several SIMD lanes in lockstep, like a GPU thread group inside one
core. Its solver takes 0.0115\,ms per round against C's 0.0183\,ms in the same session, 1.6 times faster, the
only CPU solver to beat C. It executes 163{,}000 instructions per round against C's 324{,}000, with 358 branch
mispredictions against 794. The lanes cannot simply be pointed at Algorithm~\ref{alg:reference}, because the
first derivation of each value must win. The solver therefore stores values only, so that removing
duplicates is an order-free bit OR, and finds each origin again for the answer, which is exact by
Proposition~\ref{prop:sets}. The longer half of each pair loop is spread across the lanes, which compute all four
results, their legality and the reciprocal division test (Lemma~\ref{lem:reciprocal}) together. Sorted lists
are read straight out of each subset's bitset, partner checks become single bit tests, and the sets of sizes
$n - 1$ and $n - 2$ are built but never listed, since by the counting after Lemma~\ref{lem:defer} searches
only look them up. A scalar C version of the same structure runs at 83{,}000 cycles per round against
62{,}000 for ISPC and about 95{,}000 for the C solver, so about a quarter of the gain comes from the lanes
and an eighth from a structure C could adopt. The best gang was the \emph{narrowest}: with four 64-bit lanes,
eight 32-bit lanes needed 7\% more cycles and sixteen lanes 32--41\% more. Setting a bit is a read-modify-write for each lane
slot, so idle lanes cost as much as busy ones. C's trick of sending illegal results to the value 0 also turned into
a chain of dependent updates to one word under SIMD; giving each lane its own spare word saved 19\% of ISPC's cycles
but only 1--2\% of C's.

\paragraph{What the surprises have in common.}
Each one has the same explanation underneath: what matters is how many instructions the hot loops execute
and how often they mispredict a branch, not the language's general reputation. A language does well when
it can express the design of Section~\ref{sec:cpu} directly. That means unboxed integer arrays, a population
count, and branchless selection. It does badly when its runtime inserts work into those loops: tag bits,
type tests, generic comparison, allocation zeroing or indirect calls.

\section{Sharing one round across threads}\label{sec:parallel}

The GPU solvers of Section~\ref{sec:gpu} are fast only on batches. Seven of the CPU solvers (C, Rust, Go,
Java, Julia, Erlang and Elixir) also have a separate version that splits a \emph{single} round across
threads. They use three kinds of independence in Algorithm~\ref{alg:reference}, chosen so that the answer
does not depend on which thread does what.
\begin{enumerate}[label=(\roman*)]
\item \emph{Builds within a level.} The value set of $S$ depends only on the sets of its proper subsets, so
all subsets of one size can be built at once, each by a single thread in the usual order. Removing
duplicates inside one subset stays serial, because ``first derivation found'' needs a fixed order.
\item \emph{Searches across splits.} The splits of $\Sigma(S)$ are searched by several threads, and the
earliest hit in split order wins. A shared atomic records the best position found so far, and threads skip
splits that come after it.
\item \emph{The closest pass.} Each thread finds the minimum canonical key $\kappa$ over its share of the
splits, and the minima are combined, as in the derivation-by-minimum argument used by the GPU solvers.
\end{enumerate}
Creating threads inside a solve costs more than C's entire solve, so each version keeps a worker pool alive
between solves, holding no round data. Rounds that finish within four numbers stay on one thread.

\begin{table}[ht]
\centering
\small
\begin{tabular}{@{}lrrrrrrrr@{}}
\toprule
Solver & 1 & 2 & 4 & 8 & 16 (P+HT) & 16 (P+E) & 32 & Best \\ \midrule
C       & 0.0187 & 0.0159 & 0.0142 & \textbf{0.0132} & 0.0153 & 0.0148 & 0.0236 & 1.42$\times$ \\
Julia   & 0.0208 & 0.0173 & \textbf{0.0154} & 0.0161 & 0.0181 & 0.0180 & 0.0233 & 1.35$\times$ \\
Rust    & 0.0461 & 0.0326 & 0.0253 & \textbf{0.0228} & 0.0258 & 0.0290 & 0.0416 & 2.02$\times$ \\
Java    & 0.0334 & 0.0275 & \textbf{0.0222} & 0.0266 & 0.0342 & 0.0421 & 0.0465 & 1.50$\times$ \\
Go      & 0.0296 & 0.0272 & \textbf{0.0267} & 0.0317 & 0.0429 & 0.0544 & 0.0660 & 1.11$\times$ \\
Elixir$^c$ & 0.0549 & 0.0416 & \textbf{0.0316} & 0.0318 & 0.0469 & 0.0440 & 0.0595 & 1.74$\times$ \\
Erlang$^c$ & 0.0568 & 0.0444 & \textbf{0.0330} & 0.0355 & 0.0539 & 0.0453 & 0.0761 & 1.72$\times$ \\ \bottomrule
\end{tabular}
\caption{Milliseconds per round for the parallel solvers by number of threads, best of three passes over
the 10{,}000 benchmark rounds, with every answer identical to $\operatorname{Solve}$. Up to 8 threads, each
thread has its own performance core. At 16, P+HT uses both hardware threads of the eight performance cores
and P+E adds eight efficiency cores. ``Best'' is the largest speed-up over the same code on one thread.
$^c$As in Table~\ref{tab:perf}.}\label{tab:scale}
\end{table}

Table~\ref{tab:scale} shows the result on the 32 logical CPUs of the test machine. Its findings are as follows.

\paragraph{The gain is small and stops at four to eight threads.}
The best speed-ups are between 1.1 and 2.0, and every solver is at its fastest with 4 or 8 threads. The
reason is Amdahl's law applied to Table~\ref{tab:sets}'s shape. About 46\% of rounds are settled by the
four-number subsets and never go parallel. At the five-number level there are only six subsets, so (i) can
keep at most six threads busy. Within each subset, removing duplicates is serial, and in Rust it alone
was 37\% of the time. Rounds that use all six numbers gain the most, 1.7--2.9 times at four threads, and so
do the rounds that miss, which run the closest pass (iii).

\paragraph{More threads make it slower.}
At 32 threads only Rust still beats its single thread. C is 1.26 times slower than on one thread, Erlang
1.34, Java 1.39 and Go 2.2. Each phase ends at a barrier and so waits for its slowest thread, which past 8 is
either an efficiency core, running at roughly half the speed, or a hardware thread sharing its core with
another worker. Every extra thread also adds a hand-off, scratch memory to allocate for every solve (the no-reuse
rule costs C 44\% more build work) and pressure on the shared cache, while having no work to take on.
Hardware threads and efficiency cores did about equally badly at 16 for C, Julia and Elixir. The
garbage-collected Go and Java did worse with efficiency cores, and Erlang better; we did not measure why.

\paragraph{Only spinning pools pay.}
A round lasts 20--60\,\textmu s, and a thread that has gone to sleep takes several microseconds to wake.
Handing work to a hand-written pool whose idle workers spin cost 0.3\,\textmu s in Rust and 0.6\,\textmu s in
Java, against 8.5\,\textmu s for Rust's \texttt{rayon} and 11--12\,\textmu s for Java's fork-join pool and
virtual threads. With fork-join, the parallel solver was slower than the serial one. Go's goroutines also lost to a
spinning pool. Julia's \texttt{@spawn} came closest, 7\% behind.

\paragraph{The memory manager decides the managed languages.}
At Go's default collector setting, four workers were slower than one, because each allocates a bitset of up
to about 1\,MB per round and the collector ran far more often; the pool now raises its threshold. Julia's
collector stops every thread, so its per-thread bitsets are allocated outside it, with \texttt{calloc}.
On BEAM, processes share no memory. The first parallel Erlang copied every built value list to each worker
and gained only 1.4 times. Like Elixir, it now has each worker rebuild the small subsets itself from the
six numbers and the target, pass the three-number subsets once as a shared binary, and claim tasks through
atomic counters. With that design the two BEAM languages match again, at 1.72 and 1.74 times. Both also
needed scheduler settings: idle BEAM schedulers sleep, and Elixir's spare ``dirty'' schedulers spun and took
processors from the real work, so 2 threads were slower than 1 until they were turned off.

\paragraph{A race that only 32 threads exposed.}
The first sweep never finished Elixir at 32 threads: the virtual machine spun on every CPU for three hours.
The caller read the best position of (ii) once and then waited for each split before it to be marked done.
Meanwhile one worker could lower the best, and another would then skip a split at or after the new
best without marking it, correctly, since it can no longer matter. The caller kept waiting for that split
against the stale best. Re-reading the best on every poll and stopping on reaching it removes the race.
The answer cannot change, because the winning split's result is stored before the best is lowered. The
window opens only when about 30 workers really run at once, which is why 4 or 8 threads never showed it.
The Erlang port, which copied the design, hit the same deadlock about once in 15 runs with four schedulers
before it received the same fix.

\paragraph{Summary.}
Parallelism inside one round of a single game gives at most a factor of two, and only by careful
engineering. A batch is a different matter: the CUDA solver of Section~\ref{sec:gpu} is about forty times
faster than one core, because there each round is independent and needs no barriers.

\begin{thebibliography}{9}
\bibitem{hutton} G.~Hutton. The countdown problem. \emph{Journal of Functional Programming}, 12(6):609--616, 2002.
\bibitem{repo} The Countdown solvers, benchmarks and this paper's source. \url{https://github.com/ZippyManaic/countdown}, 2026.
\end{thebibliography}

\end{document}
